\section{12.0 ZDOBYWANIE DOŚWIADCZENIA}

Ogólna zasada:

Doświadczenie (w formie punktów) jest zdobywane w walce, poprzez pokonywanie potworów.

Wielkość zdobytego doświadczenia (liczba punktów) może na końcu gry decydować o indywidualnym zwycięstwie jednego z graczy.
Ponadto doświadczenie może być użyte w celu zwiększenia poszczególnych współczynników, które charakteryzują postać.

Sposób postępowania:

Ogólna suma punktów doświadczenia, zdobyta podczas walki równa się sumie współczynników siła wszystkich zabitych potworów, pomnożonej przez 6.
Ta ogólna suma jest dzielona (w sposób, ustalony przez graczy) między pozostałych przy życiu członków oddziału.
Ułamki zaokrąglone są w dół.

Na karcie cech każdej postaci musi być prowadzony bieżący zapis sumy punktów doświadczenia, zdobytych przez daną postać od początku gry.

Jeżeli gra składa się z wielu przygód (czyli postacie wielokrotnie wchodzą i wychodzą z labiryntu) punkty doświadczenia mogą być użyte w celu zwiększenia poszczególnych współczynników, charakteryzujących postać.

Zasady szczegółowe:

[12.1] Potwory, znajdujące się pod wpływem uroku i walczące z oddziałem nie zdobywają punktów doświadczenia za zabite przez nich przeciwników – punkty te przypadają postaciom.

[12.2] Po zakończeniu gry (po wyjściu z labiryntu) postać może ,,wydać'' 100 punktów doświadczenia oraz 100 sztuk złota (lub ekwiwalent w kosztownościach), by podnieść jeden ze swoich współczynników o jeden punkt (patrz 17.5).
Wydając odpowiednią wielokrotność wymienionych sum można zwiększyć kilka współczynników, ale każdy o tylko o jeden punkt.
Tak więc zwiększeniu może ulec siła, odporność, biegłość w jednej broni, zdolności i potencjał magiczny w jednym z jego trzech członów (jeżeli postać chce zwiększyć wszystkie trzy cyfry mony musi wydać 300 punktów doświadczenia i 300 sztuk złota).

Na karcie cech należy zaznaczyć liczbę posiadanych punktów doświadczenia oraz sztuk złota, pozostałych po zapłaceniu za zwiększenie współczynników.

[12.3] Zręczność może być zwiększona podobnie jak wszystkie pozostałe cechy, jednak ,,kosztuje'' to 150 punktów doświadczenia oraz 150 sztuk złota (lub ekwiwalent w kosztownościach) – patrz 17.5
